\section{Introduction}
\label{sec:introduction}

High-frequency limit-order-book (LOB) dynamics have long been modeled with
two largely separate toolkits. The econometrics tradition favors
parametric, likelihood-based point-process models -- most prominently the
multivariate Hawkes self-exciting process -- which impose a structured,
interpretable form on order-flow clustering and cross-excitation and are
fit by unsupervised maximum likelihood on the event stream itself. The
machine-learning tradition instead favors unconstrained, supervised
function approximators -- gradient-boosted trees, in particular -- fit
directly against a forecasting target on a rich, hand-engineered feature
bank, with no likelihood story and no structural constraint on how inputs
combine. Both approaches are used, sometimes on the same data, to answer
essentially the same applied question: how likely is a given order-book
queue to deplete in the next tens to hundreds of milliseconds?

This paper asks a specific, falsifiable version of the comparison: under
what sample-size regimes, if any, does a maximum-likelihood-constrained
multivariate Hawkes process outperform an unconstrained LightGBM model at
forecasting L1 queue depletion on Binance BTCUSDT perpetual futures, when
both models are built from the \emph{same} underlying multi-scale
exponentially-weighted-moving-average (EWMA) event-count feature bank?

\subsection{The real question is generalization under constraint, not information}
\label{sec:intro-structural-note}

The comparison is deliberately not framed as ``does the Hawkes process see
something the tree cannot.'' It structurally cannot. With the same decay
grid $\beta_p \in \{0.2, 1.0, 5.0, 25.0, 100.0\}\ \mathrm{s}^{-1}$ used by
both models, the fitted multivariate Hawkes intensity
\[
  \boldsymbol{\lambda}(t) = \boldsymbol{\mu}(t) + \sum_{p=1}^{5} \mathbf{A}^{(p)} \mathbf{Z}^{(\beta_p)}(t)
\]
is, by construction, a \emph{linear} function of exactly the same
multi-scale event-count feature bank $\mathbf{Z}^{(\beta_p)}(t)$ that the
gradient-boosted-tree model ($M_1$, Section~\ref{sec:methodology}) is free
to combine nonlinearly, interact, and split on without restriction. A
model whose feature space is a strict linear subspace of another model's
own input space cannot, in the general case, exceed that other model on
any supervised objective by exploiting genuinely new information -- at
best it can match it. The interesting question this comparison can
actually answer is therefore not an information question but a
\emph{generalization} question: does the regularization implicit in an
unsupervised, likelihood-constrained parametric fit -- effectively, a
structured dimensionality reduction of the same feature space down to a
handful of branching-ratio parameters per event-type pair -- generalize
better than letting a flexible nonparametric learner fit that same space
directly, particularly in small-sample, regime-shifted, or
cross-instrument-transfer conditions where the tree has less to learn
from and a strong parametric prior might be expected to help most?

\subsection{Falsification-first framing}
\label{sec:intro-falsification}

This project is explicitly designed to be falsification-first rather than
confirmation-seeking. A single Primary Confirmatory Endpoint is
pre-registered before any result is examined: on the 500\,ms-horizon,
50\%-of-spread adverse-price-shift target, a treatment model must beat
$M_1$ by $\Delta\PRAUC \ge +0.005$ at $\alpha = 0.05$ under a formal,
walk-forward significance test (the rolling Giacomini--White protocol of
Section~\ref{sec:methodology}) for the parametric-compression hypothesis
to be considered supported. Every other comparison in this paper --
the full sample-size learning curve, the 18-specification secondary grid,
the cross-asset transfer ladder, the operational lead-time metric -- is
reported as exploratory evidence, not as a confirmatory test, and is
labeled as such throughout. Where results were surprising or ran counter
to a standing project hypothesis -- most notably, a documented
expectation that untuned LightGBM hyperparameters were flattering $M_1$,
later tested and found backwards (Section~\ref{sec:results-tuning}) -- the
finding is reported plainly rather than smoothed over, consistent with
the project's own stated philosophy that the goal is not to prove the
Hawkes process adds value but to test, honestly, whether it does.

\subsection{Why this matters}
\label{sec:intro-why}

Two practical audiences care about the answer. First, for real-time
microstructure forecasting infrastructure, the choice between a
parametric point-process model and a free tree model is a choice about
sample-efficiency, interpretability, and cross-instrument transferability
versus raw predictive power -- a Hawkes process needs far fewer bytes to
describe, is cheaper to fit incrementally, and, as
Section~\ref{sec:results-transfer} shows, transfers its fitted structure
across instruments in a way a black-box tree model has no native
mechanism to do at all. If the parametric model's generalization
advantage in the small-sample or cross-instrument regime is real, that
trade-off favors the parametric approach in exactly the conditions where
a new instrument or a cold-start deployment has little training data of
its own. Second, for the broader methodological question of when
structured, likelihood-based inductive bias helps versus hurts relative
to unconstrained supervised learning on the same feature space, L1 queue
depletion is a clean, data-rich testbed: the ground-truth event stream is
large (in excess of one billion typed events across the primary window),
the two models share an identical feature basis by design, and the
forecasting target is well defined and mechanically close to the
underlying point process. A clear empirical answer in this setting speaks
to a question broader than any one exchange or instrument.

The remainder of the paper is organized as follows.
Section~\ref{sec:related-work} situates this work relative to prior
literature on Hawkes processes in market microstructure, LOB
queue-depletion modeling, and machine-learning approaches to
high-frequency forecasting. Section~\ref{sec:data} describes the data
pipeline, the 38-day primary window, the 10-event microstructure
taxonomy, and known data-quality constraints. Section~\ref{sec:methodology}
gives the full model ladder ($M_0$--$M_3$), the Hawkes MLE formulation,
the LightGBM setup, and the statistical testing protocols.
Section~\ref{sec:results} reports the empirical results: the fixed-test
learning curve, the rolling Giacomini--White test, the calibration-budget
robustness sweep, the $M_3$ build and its ablation, the 18-specification
Romano--Wolf grid, the cross-asset transfer ladder, the per-$N$
hyperparameter-tuning finding, and the operational lead-time metric.
Section~\ref{sec:limitations} discusses scope limitations, disclosed
simplifications and their resolutions, and reproducibility caveats.
Section~\ref{sec:conclusion} concludes.
